\section{Mechanism}
\label{sec:mechanism}

The compiled channel is a hypernetwork that amortizes into one forward pass what LoRA
obtains by gradient descent: a trained write function maps a structured belief
representation, through the frozen host's own hidden states, to per-layer conditioning
parameters, in the manner of AdaLN- and FiLM-style conditional normalization
\citep{perez1970film,xu2019adain} and of the hypernetwork-adapter lineage established by
GenerativeAdapter \citep{chen2025generativeadapter} and Text-to-LoRA
\citep{charakorn2025texttolora}.

\subsection{The belief tree}
\label{sec:mechanism-belieftree}

Conditioning content is stored as a \emph{belief tree}: typed hierarchical nodes
\(\{\text{statement}, \text{type} \in \{\text{claim}, \text{argument}, \text{evidence},
\text{experience}, \text{strategy}\}, \text{credence}, \text{edges}
\{\text{supports}, \text{contradicts}\}\}\). Trees are authored by hand for evaluation
topics and produced by the reflection loop for the update-path experiments
(Section~\ref{sec:discrimination}). An early design used numeric credence values on
every node; we found no measurable benefit over a coarse categorical scheme and report
this as a negative result rather than carry the added complexity forward.

\subsection{The modulation}
\label{sec:mechanism-modulation}

The write function conditions query and value projections at a stride-3 subset of
layers (11 of 32), rank 16. For a target layer with base query weight \(W_Q\), the
modulated projection is
\[
  W_Q' = W_Q + \Delta_Q(\phi), \qquad \Delta_Q(\phi) = U_Q(\phi)\,V_Q(\phi)^\top,
\]
with \(U_Q(\phi) \in \mathbb{R}^{d \times 16}\), \(V_Q(\phi) \in \mathbb{R}^{d \times
16}\) produced by the write function from the pooled belief-tree state \(\phi\); the
value projection \(W_V'\) is constructed identically. The compilation is
query-\emph{agnostic} --- \(\Delta_Q\), \(\Delta_V\) do not depend on the current
prompt --- but its effect is query-\emph{dependent}, because the low-rank update acts
on the live hidden state produced by whatever the model is currently attending to.

\subsection{The write function}
\label{sec:mechanism-write}

Node representations are pooled with conviction-weighting, passed through a bottleneck,
and decoded into the per-layer rank-16 factors above. We correct a timing claim from an
earlier draft of this work: the previously reported $\sim$20 second compile latency
belongs to an unrelated agentic retrieval prototype, not to this write function, which
compiles in well under one second; end-to-end reflection $\to$ edit $\to$ recompile
measures 2.5--3.0 seconds.

\textbf{We define ``compile'' at first use, and it is lossy.} A trained write function
translates the belief tree into per-layer low-rank tensors, fixed until the store
changes, yet query-dependent in effect because the update acts on the live hidden state.
We reject ``precomputed,'' which suggests caching a value that could in principle be
looked up; this is a learned, lossy translation, and the loss is what
Section~\ref{sec:boundary} characterizes. We also flag, to preempt confusion with an
unrelated sense of the word, that ``compile'' here has no relation to \texttt{torch.compile}
or graph compilation.

\subsection{The curated slot}
\label{sec:mechanism-slot}

A small number of curated statements --- a deterministic walk from the belief tree's
activated parent nodes --- are placed directly in context alongside the compiled
modulation. This slot carries what the bottleneck cannot: specific, low-frequency
content such as named entities, exact figures, or constrained action strings
(Section~\ref{sec:boundary}). Condition \condF{} throughout this paper denotes compiled
modulation plus this curated slot; \condC{} denotes compiled modulation with the slot
empty (zero belief-content context tokens).

\begin{figure}[t]
\centering
\fbox{\parbox{0.9\textwidth}{\centering \fixme{Figure 1: compile pass (belief tree
$\to$ pooled state $\to$ write function $\to$ per-layer $\Delta_Q,\Delta_V$) and
generation pass (modulated forward pass with the curated slot in context), side by
side, readable without the caption. Insert
\texttt{/home/dgonier/experiments/figures/} asset once
\texttt{paper\_figures\_dispositional.py} emits it.}}}
\caption{The compile pass and the generation pass. Compilation happens once per belief
change (2.5--3.0s end-to-end); generation applies the resulting per-layer tensors on
every subsequent forward pass at no additional context cost.}
\label{fig:mechanism}
\end{figure}
